\documentclass{subfiles}
\begin{document}
    In the previous section we discussed Canonical Huffman encoding, 
        which solves the issue of transmitting the tree, 
        but not that of scanning the text twice.
        In this section we shall present a variant of Huffman encoding that 
        solves this issue.

    Proposed independently by Fallar in 1973 and Gallager in 1978,
        later improved by Knut, the Adaptive Huffman encoding
        solves the issue of repeated scans of the text own of the classic approach.
        The core idea is simple: build a optimal Huffman tree for the text 
        process thus far, and reorganize the tree when necessary.

    More in details, we can summarize the procedure as follow:
    \begin{enumerate}
        \item Let \(T\) be an empty tree, or let it contain a special symbol 
            NYT.
        \item While scanning the text, for each symbol:
            \begin{itemize}
                \item If it has already occurred, use the available codeword to encode it;
                \item If not, encode the NYT symbol, then add the symbol as a new leaf and encode it.
            \end{itemize}
        \item After encoding a symbol, update \(T\) as follow:
                \begin{itemize}
                    \item Increment the frequency of the symbol and all its ancestors;
                    \item Reorganize the tree to maintain the Huffman properties.
                \end{itemize}
    \end{enumerate}

    It can be proved that, in general,
        the adaptive approach and the classic one have the same performances.
\end{document}
